\chapter{Differentiable tracking}\label{ch:differentiable}

\section{Why gradients}
Calibrating a model on machine data means fitting a few parameters (field scale, dee voltages, phases, the initial beam)
to measured quantities (probe currents, phase, extracted energy).
With gradients of the tracked observables with respect to these parameters, the fit can use gradient-based optimisers.
With $10^5$ particles, the forward model can also be fast enough for many evaluations.
Differentiable beam dynamics codes exist for linear accelerators and transfer lines \cite{kaiser2024}.
Cases E and E2 ask whether the LIONS core can join them.

\section{Reverse-mode automatic differentiation}
A tracking run is a composition $X_N = \Phi_N\circ\cdots\circ\Phi_1(X_0;\lambda)$.
Reverse mode computes $\partial f(X_N)/\partial\lambda$ for a scalar $f$ in one backward sweep, at a cost of a few forward
passes, whatever the number of parameters.
The price is memory: the backward sweep needs the intermediate states of all $N$ steps.
Checkpointing stores only every $\sqrt N$-th state and recomputes each segment during the backward sweep.
In principle memory then grows as $O(\sqrt N)$ for one extra forward pass \cite{chen2016}; recursive schemes reach
logarithmic growth \cite{griewank2000}.

\section{Case E: the port and its gradients}
The LIONS bilinear field and Boris push were ported to PyTorch \cite{paszke2019} in float64.
The port reproduces LIONS positions and momenta at every step of one turn within \EPortMax{} relative.
Gradients of $r$, $p_r$ and the phase with respect to the map scale $s$, $r_0$ and $p_{r0}$ were compared with central finite
differences (\cref{tab:Egrad}).

\begin{table}[h]
\centering\small
\caption{Case E: maximum relative difference between autograd and central finite differences at the registered steps
($10^{-6}$ relative for $s$, \SI{1e-7}{m} for $r_0$, $10^{-7}$ for $p_{r0}$).}
\label{tab:Egrad}
\input{gen/tab_e_grad.tex}
\end{table}

\begin{itemize}
\item B0, the smooth field, passes the $10^{-6}$ criterion.
\item A, the uniform field, fails on \EAFailCount{} entries.
The largest are $\partial\phi/\partial p_{r0}$ at both step counts and $\partial\phi/\partial r_0$ at 3600 steps.
These derivatives are small: in a uniform field, at fixed $|\vect p|$, the revolution period does not depend on $r_0$ or
$p_{r0}$, while the finite-difference reference carries a rounding error that grows with the phase (about $2\pi$ per
turn): there the finite difference is the noisy side.
The other failures are at most \EAFailOtherMax{} and were not analysed.
\item B, the bilinear 4-sector map, is the instructive case (\cref{fig:Efd}).
At 3600 steps the disagreement reaches \ESweepLargeH{} at $h = 10^{-4}$, falls to \ESweepSmallH{} at $h = 10^{-8}$, and
rises again at $h = 10^{-9}$ (\ESweepTinyH) from rounding.
(The summary of this sweep and of the case A failures in \nolinkurl{benchmarks/DESIGN.md}, section 19, does not match
the table; the numbers here are from the table.)
Example: $\partial r/\partial s = \EExampleAutograd$ from autograd, against \EExampleFD{} from finite differences at
$h = 10^{-6}$.
\end{itemize}

\begin{figure}[t]
\centering
\includegraphics[width=0.6\textwidth]{case_e_fd}
\caption{Case E: worst disagreement between autograd and central finite differences against the finite-difference step,
at 3600 steps (10 turns).
The bilinear map (case B) agrees only near $h=10^{-8}$; on the smooth field (B0) the disagreement stays at the
finite-difference error level (truncation at large $h$, rounding at small $h$).}
\label{fig:Efd}
\end{figure}

\begin{reading}
Autograd returns the derivative of the discrete map at the point.
On a bilinear map the discrete map is only piecewise $C^1$: its derivative jumps whenever a step point crosses a cell
edge.
A finite difference with a step larger than the spacing of these jumps averages over them.
For an optimiser this means the local gradient and the gradient over a finite parameter change differ: after ten turns,
by up to \ESweepLargeH{} for a relative change of $10^{-4}$ of the map scale $s$.
A $C^1$ field (F1) removes the jumps; whether it closes the gap for gradients has not been measured.
\end{reading}

\paragraph{Memory and time.}
Without checkpointing the autograd graph costs about \EKBPerStep~kB per step for one particle (from the $10^3$ and $10^4$
step runs), which \nolinkurl{benchmarks/DESIGN.md} (section 19) attributes to about 60 small operations per step.
The measurements (\cref{tab:Ecost}) do not follow the textbook picture.
Checkpointing reduces memory only with $10^3$ particles; with one particle it uses more.
It costs \ECkptTimeMin{} to \ECkptTimeMax{} times the time, not about two.
At $10^5$ steps it still needs several GB above the imports, far more than $\sqrt N$ stored states.
The cause was not investigated.

\begin{table}[h]
\centering\small
\caption{Case E: forward plus backward, eager PyTorch on the M3 CPU (peak RSS includes the Python and PyTorch imports).
``skipped'': extrapolated memory above half the RAM.}
\label{tab:Ecost}
\input{gen/tab_e_cost.tex}
\end{table}

\section{Case E2: throughput}
\textbf{Question} (registered before running): can a compiled PyTorch step beat plain LIONS throughput by 10 times?
This is the go criterion for building a differentiable backend.

\textbf{Result} (\cref{fig:E2}, \cref{tab:E2}).
Compiled float64 reaches \EtwoRatioH{} times plain LIONS at $10^5$ particles on an H100 (\EtwoBestH{} particle-steps per
second), \EtwoRatioA{} times on an A100, and \EtwoRatioM{} times on the M3 CPU.
At $10^3$ particles no device reaches 10.
float32 gives no speed gain on the GPUs and loses accuracy: up to \EtwoGpuFloatThirtyTwoR{} relative in $r$ after ten
turns.
GPU float64 agrees with CPU float64 within \EtwoGpuFloatSixtyFourR.

\begin{table}[h]
\centering\small
\caption{Case E2: forward throughput relative to plain LIONS on the same node, and wall time of compiled float64
forward plus backward at $10^3$ particles (3600 steps).}
\label{tab:E2}
\input{gen/tab_e2.tex}
\end{table}

\begin{figure}[t]
\centering
\includegraphics[width=0.65\textwidth]{case_e2}
\caption{Case E2: forward particle-steps per second against particle count, best float64 variant per device
(MPS: float32 only).
Dashed: plain LIONS on the CPU of the same node.}
\label{fig:E2}
\end{figure}

\begin{reading}
The go criterion is met at $10^5$ particles, on the GPUs and on the M3 CPU.
On the A100 the wall time changes from \EtwoAWallOne{}~s to \EtwoAWallMax{}~s between 1 and $10^5$ particles, so the run is launch-bound and $10^5$ is not the
ceiling.
Small-batch gradients are faster on the CPU.
\end{reading}

\section{Before trusting a calibrated model}
A model fitted on machine data inherits every omission of the simulator.
LIONS does not model space charge, a full 3D field map, misalignments of the dees or the central-region electrodes,
field drifts and hysteresis, or RF amplitude and phase noise.
A parameter fit will absorb these effects into the parameters it has.
Fitted values are therefore effective parameters until they are checked on held-out operating conditions.
